\section{Model}\label{sec:model}

This section develops a tractable model of startup subsidies, buyout exit, and merger review.
The key ingredients are a local entry subsidy, a post-entry choice between an independence-oriented and a buyout-oriented project, and a reduced-form merger approval probability that governs the attractiveness of acquisition as an exit route.

\subsection{Environment and timing}

Consider a unit measure of independent local markets indexed by \(m \in [0,1]\).
In each market there is one incumbent firm and one potential startup project.
Markets are otherwise identical, except that potential entrants differ in their fixed entry costs.
A local government offers a uniform entry subsidy \(s \ge 0\) to any entrepreneur who enters.

In each market, the timing is as follows:

\begin{enumerate}
    \item The local government chooses a uniform entry subsidy \(s\).
    \item The potential entrepreneur in market \(m\) observes an idiosyncratic fixed cost \(f_m\) and decides whether to enter.
    \item Conditional on entry, the entrepreneur chooses a project type \(j \in \{I,B\}\) and an effort level \(x \in [0,1]\).
    \item If the project succeeds, the incumbent and the entrant bargain over a potential merger.
    \item If a merger is attempted, it is approved with probability \(q \in [0,1]\).
    \item If the merger is approved, the merged firm operates in the product market; otherwise the incumbent and the entrant compete as independent firms.
\end{enumerate}

We interpret \(q\) as a reduced-form measure of merger permissiveness.
A higher \(q\) corresponds to a more permissive merger environment, or more generally to an environment in which buyout exit is easier to realize.
The baseline keeps this shift project-invariant in order to focus on project choice and welfare accounting in a one-dimensional way; Appendix \ref{app:typespecificq} briefly notes how the mechanism extends when approval or exit realizability differs across project types.
To avoid confusion, \(q\) is the approval or exit-realizability parameter, whereas \(q_I\) and \(q_E\) below denote product-market quantities.

The entry cost \(f_m\) is independently drawn from the uniform distribution on \([0,1]\).
The uniform specification is used for closed-form results in the main text; the appendices show how the analysis extends to a general distribution \(G(f)\).

\subsection{Product-market payoffs}

If entry succeeds, the product-market primitives depend on the entrepreneur's project choice \(j \in \{I,B\}\).
Type \(I\) denotes an \emph{independence-oriented} project: it is relatively strong as a stand-alone business.
Type \(B\) denotes a \emph{buyout-oriented} project: it is relatively attractive as an acquisition target.
These two project types should be read as reduced-form project orientations rather than as a fully dynamic innovation-race model.
The primitive differences \((a_j,\sigma_j,m_j)\) capture, respectively, stand-alone marketability, product overlap with the incumbent, and acquisition-specific integration value or rents.

For a successful entrant of type \(j\), inverse demands are given by
\begin{align}
    p_I &= 1 - q_I - \sigma_j q_E, \\
    p_E &= a_j - q_E - \sigma_j q_I,
\end{align}
where \(q_I\) and \(q_E\) denote the incumbent's and entrant's quantities, respectively,
\(a_j > 0\) is the entrant's stand-alone demand shifter, and \(\sigma_j \in (0,1)\) captures substitutability between the incumbent's and the entrant's products.
Marginal costs are normalized to zero.

We impose parameter restrictions ensuring interior product-market solutions under both independent competition and common ownership:
\begin{equation}\label{eq:interiorpm}
    \sigma_j < a_j < \frac{1}{\sigma_j},
    \qquad j \in \{I,B\}.
\end{equation}

\paragraph{Independent competition.}
If no merger takes place, the incumbent and the entrant compete \`a la Cournot.
The incumbent solves
\begin{equation}
    \max_{q_I \ge 0}\; (1-q_I-\sigma_j q_E)q_I,
\end{equation}
and the entrant solves
\begin{equation}
    \max_{q_E \ge 0}\; (a_j-q_E-\sigma_j q_I)q_E.
\end{equation}
The unique interior Nash equilibrium is
\begin{align}
    q_I^D(j) &= \frac{2-a_j\sigma_j}{4-\sigma_j^2}, \\
    q_E^D(j) &= \frac{2a_j-\sigma_j}{4-\sigma_j^2}.
\end{align}
At this equilibrium, prices satisfy
\begin{equation}
    p_I^D(j)=q_I^D(j),
    \qquad
    p_E^D(j)=q_E^D(j),
\end{equation}
so profits are
\begin{align}
    \pi_I^D(j) &= \left(\frac{2-a_j\sigma_j}{4-\sigma_j^2}\right)^2, \\
    \pi_E^D(j) &= \left(\frac{2a_j-\sigma_j}{4-\sigma_j^2}\right)^2.
\end{align}

Consumer surplus under independent ownership is
\begin{equation}\label{eq:csd}
    C^D(j)
    =
    \frac{1}{2}
    \left[
        \bigl(q_I^D(j)\bigr)^2
        + \bigl(q_E^D(j)\bigr)^2
        + 2\sigma_j q_I^D(j)q_E^D(j)
    \right].
\end{equation}

\paragraph{Approved merger.}
If the entrant is acquired and the merger is approved, the merged firm chooses \((q_I,q_E)\) to maximize joint operating profit.
In addition, common ownership generates a type-specific synergy \(m_j \ge 0\), which captures the private value of integrating the entrant's technology, customer base, team, or complementary assets into the incumbent.

The merged firm solves
\begin{equation}
    \max_{q_I,q_E \ge 0}\;
    (1-q_I-\sigma_j q_E)q_I
    +
    (a_j-q_E-\sigma_j q_I)q_E
    +
    m_j.
\end{equation}
The unique interior solution is
\begin{align}
    q_I^M(j) &= \frac{1-a_j\sigma_j}{2(1-\sigma_j^2)}, \\
    q_E^M(j) &= \frac{a_j-\sigma_j}{2(1-\sigma_j^2)}.
\end{align}
The corresponding operating profit is
\begin{equation}
    \Pi_{\mathrm{op}}^M(j)
    =
    \frac{a_j^2 - 2a_j\sigma_j + 1}{4(1-\sigma_j^2)},
\end{equation}
so total merged profit is
\begin{equation}
    \Pi^M(j) = \Pi_{\mathrm{op}}^M(j) + m_j.
\end{equation}
Consumer surplus under common ownership is
\begin{equation}\label{eq:csm}
    C^M(j)
    =
    \frac{1}{2}
    \left[
        \bigl(q_I^M(j)\bigr)^2
        + \bigl(q_E^M(j)\bigr)^2
        + 2\sigma_j q_I^M(j)q_E^M(j)
    \right].
\end{equation}

Define the private merger surplus for project \(j\) as
\begin{equation}\label{eq:Deltaj}
    \Delta_j \equiv \Pi^M(j) - \pi_I^D(j) - \pi_E^D(j).
\end{equation}

The two project types differ in their stand-alone and acquisition properties.
Throughout, we assume that type \(I\) is more attractive under independent ownership, whereas type \(B\) generates a larger private gain from acquisition.

\begin{assumption}\label{ass:types}
The product-market primitives satisfy
\begin{equation}
    \pi_E^D(I) > \pi_E^D(B)
    \qquad \text{and} \qquad
    \Delta_B > \Delta_I > 0.
\end{equation}
\end{assumption}

Assumption \ref{ass:types} captures the intended distinction between the two project types.
A sufficient primitive configuration is that type \(I\) has a higher stand-alone demand shifter, whereas type \(B\) features a larger synergy term \(m_j\) and/or higher substitutability \(\sigma_j\), making it more valuable under common ownership.
In that sense, the model's notion of entrepreneurial direction is the direction of project orientation between stand-alone competition and buyout exit.

\subsection{Project choice, effort, and entry}

After entering, the entrepreneur chooses a project type \(j \in \{I,B\}\) and an effort level \(x \in [0,1]\).
Effort is interpreted as the probability of technical and commercial success.
Its cost is
\begin{equation}
    c(x) = \frac{x^2}{2}.
\end{equation}

If the project succeeds, the entrepreneur obtains a project-specific continuation payoff that depends on whether a merger is attempted and approved.
Let \(R_j(q)\) denote the entrepreneur's expected continuation payoff conditional on success and project choice \(j\).
Given \(R_j(q)\), the entrepreneur solves
\begin{equation}
    \max_{x \in [0,1]} \; xR_j(q) - \frac{x^2}{2}.
\end{equation}

Hence, before success is realized, the entrepreneur's expected payoff from entering with project \(j\) and fixed cost \(f_m\) is
\begin{equation}
    U_j(f_m;s,q)
    =
    s - f_m
    + \max_{x \in [0,1]}
      \left\{
        xR_j(q) - \frac{x^2}{2}
      \right\}.
\end{equation}

The entrepreneur chooses the project type that maximizes \(U_j(f_m;s,q)\), and then decides whether to enter.
Thus entry occurs if and only if
\begin{equation}
    \max_{j \in \{I,B\}} U_j(f_m;s,q) \ge 0.
\end{equation}

We restrict attention to parameter values such that the optimal effort choice is interior in the main text.
Under this restriction, the entrepreneur's effort and entry decisions admit closed-form solutions characterized in the next section.

\subsection{Acquisition bargaining and merger review}

If a type-\(j\) project succeeds, the incumbent and the entrant may attempt a merger.
Filing and implementing the transaction requires a fixed cost \(k>0\), interpreted as the legal, administrative, and organizational cost of merger review and consummation.

The sequence is as follows.
After success is realized, the incumbent and the entrant observe project type \(j\) and bargain over the merger.
If they attempt a merger, approval occurs with probability \(q\); if approval is denied, the firms remain independent and compete in the product market.

We assume Nash bargaining over the \emph{expected} merger surplus, with the entrepreneur receiving bargaining weight \(\beta \in (0,1)\).
The no-merger outside options are \(\pi_I^D(j)\) and \(\pi_E^D(j)\).
If the firms attempt a merger, the expected joint payoff net of filing cost is
\begin{equation}
    q\Pi^M(j) + (1-q)\bigl(\pi_I^D(j)+\pi_E^D(j)\bigr) - k,
\end{equation}
whereas the outside option from not filing is
\begin{equation}
    \pi_I^D(j)+\pi_E^D(j).
\end{equation}
Hence the incremental expected gain from attempting the merger is
\begin{equation}
    q\Delta_j - k.
\end{equation}

The firms attempt a merger if and only if
\begin{equation}\label{eq:attemptcondition}
    q\Delta_j \ge k.
\end{equation}
Under Nash bargaining, the entrepreneur's expected continuation payoff conditional on success is therefore
\begin{equation}\label{eq:Rj}
    R_j(q)
    =
    \pi_E^D(j)
    +
    \beta \bigl[q\Delta_j-k\bigr]_+,
\end{equation}
where \([z]_+ \equiv \max\{z,0\}\).

Equation \eqref{eq:Rj} is central.
A higher approval probability \(q\) raises the attractiveness of project types with large merger surplus.
As a result, merger policy affects not only the quantity of entry, but also the direction of entrepreneurial innovation.

\subsection{Policy-relevant welfare ingredients}

The baseline policy instrument is the local entry subsidy \(s\). The government takes the merger approval probability \(q\) as given and chooses \(s\) before entry takes place.

For policy analysis, social value is built from consumer surplus, profit components, and non-pecuniary spillovers under either independent ownership or approved acquisition. Section \ref{sec:subsidy} first aggregates these primitives into a benchmark social value that depends only on product-market primitives and merger-attempt behavior. Section \ref{sec:decentralized} then adds local and national incidence weights on top of that benchmark. The crucial point is that merger attempt is endogenous: the same condition \eqref{eq:attemptcondition} that governs private continuation values must also govern welfare accounting. If \(q\Delta_j<k\), no merger is attempted and successful entry yields the independent-ownership value. If \(q\Delta_j \ge k\), a merger is attempted, approval occurs with probability \(q\), and the resource cost of attempt is incurred regardless of approval.

The next sections use this behaviorally consistent accounting to derive the primitive benchmark subsidy, then compare planner-specific local and national incidence extensions around that benchmark.
